\section{Comparison with Binius64}
\label{sec:cmp}

We close by positioning the construction against \binius{} \cite{binius64repo},
the SHA-256 prover of the Binius line \cite{binius,fribinius} and the closest
system in the literature. The two are cousins: this work imports \binius{}'s
per-step modular-addition carry identity (\Cref{sec:add}) and its single-$\AND$
form of $\Ch$, and both work in characteristic two precisely to avoid the
prime-field witness blow-up that bit-decomposing SHA-256 incurs. They differ in
\emph{one} structural choice --- \binius{} arithmetizes SHA-256 as an unrolled,
non-deterministic $64$-bit \emph{word circuit}, this work as a \emph{uniform} AIR
over the cell ring --- which, with two further independent choices of commitment
and reduction field, accounts cleanly for the numbers a user observes.

\subsection{Three orthogonal design axes}
\label{sec:axes}

The systems make independent choices on three axes, and each axis decides one
observable outcome (\Cref{tab:axes}).

\begin{table}[h]
\centering
\renewcommand{\arraystretch}{1.25}
\begin{tabularx}{\textwidth}{@{}lXXl@{}}
\toprule
\textbf{Axis} & \textbf{\binius{}} & \textbf{This work} & \textbf{Decides} \\
\midrule
Arithmetization &
  unrolled word \emph{circuit} (arbitrary wiring) &
  uniform \emph{AIR} over $\Ftwo[X]$ (structured wiring) &
  \textbf{verifier} \\
Commitment / proximity &
  Reed--Solomon $+$ recursive FRI/BaseFold &
  linear $\Ftwo$-RAA $+$ flat Brakedown &
  \textbf{proof size} \\
Reduction field &
  sumcheck tensors in $\mathrm{GF}(2^{128})$ &
  discharge kept in $\Ftwo/\mathrm{GF}(2^8)$ &
  \textbf{memory / large $N$} \\
\bottomrule
\end{tabularx}
\caption{The three design axes. Each system picks a different point on each, and
each axis decides one headline outcome. No single axis decides the prover, which
is dominated by content both systems share.}
\label{tab:axes}
\end{table}

\emph{Arithmetization decides the verifier.} A verifier is succinct exactly when
the constraint system's wiring --- the map from each constraint to the witness
positions it reads --- is succinctly evaluable at a random point. The uniform AIR
has \emph{structured} wiring: one fixed transition relation, an $O(\log N)$
formula, the verifier never reading per-row structure. \binius{}'s word circuit
has \emph{arbitrary} wiring: which committed word feeds which gate is data with no
closed form, so the verifier sums over all $N$ gates in the clear --- its $O(N)$
verifier, and by the same token its $O(N)$ preprocessing against this work's
$O(1)$ (\Cref{app:meas}). It is the cross-entry \emph{wiring}, not the entry-wise
shifts, that is linear: the entry-wise shift is succinct in both systems.

\emph{Commitment decides proof size.} Independently, this work buys a
linear-time $\Ftwo$-RAA encoder and pays with a \emph{flat} proof --- a Brakedown
test opening $t\approx 987$ full columns, growing as $t\cdot\sqrt N$ with a large
floor --- while \binius{} buys a \emph{small, recursive} FRI proof and pays with
an NTT encoder. A WHIR-style recursive proximity test would close this gap on
this side (\Cref{rem:blaze}).

\emph{Reduction field decides memory and scaling.} \binius{}'s BitAnd and Shift
reductions build $\mathrm{GF}(2^{128})$ tensors that spill cache super-linearly;
this work keeps the witness bit-packed and runs the Hadamard discharge in the
base field $\Ftwo/\mathrm{GF}(2^8)$, lifting to $\K$ only at the readings ---
roughly an order of magnitude fewer bytes streamed and about half the footprint
at the largest matched instance.

\emph{No single axis decides the prover.} The prover is dominated by work both
systems must do --- discharge the same $\AND$ content, encode the witness --- and
the word-level (one packed lane) versus bit-sliced ($\times\wbit$) asymmetry in
how they reach it roughly cancels at the current optimisation level. The prover
is thus a \emph{net} of large offsetting per-phase gaps (\Cref{app:meas}) rather
than a single dominant difference.

\subsection{Measured comparison}
\label{sec:results}

All figures are on one Apple~M4, non-zero-knowledge, multithreaded, CPU, at the
matched commitment rate $1/4$. Absolute times are thermally variable, so we read
\emph{ratios} and per-phase \emph{shapes}.

\begin{table}[h]
\centering
\renewcommand{\arraystretch}{1.2}
\begin{tabular}{@{}llrrl@{}}
\toprule
& & \binius{} & This work & winner \\
\midrule
\multirow{3}{*}{$\sim$$1000$ comp}
  & Prove   & $84$\,ms    & $75$\,ms    & this work $1.1\times$ \\
  & Verify  & $14$\,ms    & $12$\,ms    & this work $1.2\times$ \\
  & Proof   & $249$\,KiB  & $1.12$\,MiB & \binius{} $4.6\times$ \\
\midrule
\multirow{3}{*}{$\sim$$8000$ comp}
  & Prove   & $1.22$\,s   & $0.58$\,s   & this work $2.1\times$ \\
  & Verify  & $493$\,ms   & $45$\,ms    & this work $11\times$ \\
  & Proof   & $299$\,KiB  & $\sim$$2.0$\,MiB & \binius{} $\sim$$7\times$ \\
\bottomrule
\end{tabular}
\caption{Matched rate $1/4$, CPU, at two scales (full sweeps in \Cref{app:meas}).
At moderate $N$ the systems are close on every axis but proof size; the content
is how they \emph{scale}: this work's prover is linear and \binius{}'s
super-linear, so the prover gap opens to $\sim$$2\times$ by $8000$ compressions
and then becomes unbounded (\binius{}'s $O(N)$ circuit build walls near
$16{,}000$, where this work still proves $123{,}366$ via swap). The verifier and
proof move oppositely: $\sqrt N$ versus $O(N)$ verifier, flat versus recursive
proof.}
\label{tab:headline}
\end{table}

\subsection{The soundness asymmetry}
\label{sec:cmp-soundness}

A fair comparison must record that the two are not at equal soundness status.
\binius{} is fully sound at its $96$-bit target. This work is unconditionally
sound for the $\AND$ relations and the entire linear part, but its
modular-addition carries are presently \emph{trusted}: $12$ of the $14$ deployed
Hadamards have honest-prover-supplied $\psiz$ parents (\Cref{rem:adder-trusted}).
The recommended fix replaces those $12$ carry Hadamards with one grand-product
lookup against a public addition table \cite{lasso} --- sound by construction,
and simultaneously removing the $\times\wbit$ bit-slice discharge cost for
additions. A characteristic-two subtlety is essential: the lookup must use a
\emph{multiplicative} grand product $\prod_i(\gamma-a_i) = \prod_t(\gamma-t)^{m_t}$,
not an additive log-derivative, which is unsound over an $\Ftwo$-extension because
$1/(\gamma-a)+1/(\gamma-a) = 0$ lets even multiplicities cancel. Because closing
the gap costs prover time not yet paid, the moderate-$N$ edge includes a small
soundness discount; the asymptotic (scaling) win is the more robust claim.

\subsection{When to use which}
\label{sec:guidance-brief}

The three axes are \emph{hash-agnostic} --- they concern the arithmetization, the
proximity test, and the reduction field, not SHA-256 specifically --- so they
yield workload guidance. Hash-dominated work \emph{at scale} favours this work on
prover throughput, memory, and verifier cost together; a hard \emph{small
standalone proof} favours \binius{} (or a recursion layer wrapped over this work,
whose $\sqrt N$ verifier is the cheaper inner argument to wrap); \emph{heterogeneous}
statements pairing a hash with big-integer arithmetic favour this work's
multi-ring composition in a unified limb domain; and at \emph{moderate} scale the
systems are close enough that the binding output constraint decides.

\section{Conclusion}
\label{sec:conclusion}

This work and \binius{} agree on the hard decision --- prove SHA-256 in
characteristic two, refusing the prime-field witness blow-up --- and disagree on
three orthogonal downstream choices, each deciding one outcome: uniform AIR
versus word circuit decides the verifier (and $O(1)$ versus $O(N)$
preprocessing); flat versus recursive proximity decides the proof; the
base-field versus $\mathrm{GF}(2^{128})$ reduction decides memory and scaling.
The prover, dominated by shared content, is the closest of the four: at the
matched rate on one machine, this work leads by a touch at moderate $N$, widening
to $\sim$$2\times$ and then unbounded as \binius{}'s $O(N)$ build walls; its
verifier runs from near-parity to past $11\times$ ($\sqrt N$ versus $O(N)$) on
about half the memory, while \binius{}'s recursive proof stays $4$--$7\times$
smaller and it is fully sound where these adders are trusted.

Two observations resist the usual summaries. First, \binius{}'s linear verifier
is caused by the cross-entry \emph{wiring}, not the shifts: the entry-wise shift
is succinct in both systems, and the wiring is succinct only under a uniform
arithmetization. Second, the levers each system needs point toward the other ---
this work's proof wants FRI-style recursion, \binius{}'s prover wants the free
shift of a cell-ring AIR --- the convergent endpoint being a binary-field STARK
with an $\Ftwo[X]$ cell ring, a uniform AIR, a base-field discharge, a sound
lookup adder, and a recursive proximity test.
